\chapter{Introduction}
\label{ch:introduction}

\section{Problem and construct boundary}

Recorded mock interviews contain language, vocal delivery and visual behaviour that can be used
to predict how observers score a participant. Such prediction is scientifically useful only when
the target is named accurately. The first study predicts the MIT Interview corpus's aggregate
\emph{RecommendHiring} rating. That quantity is a crowd-perceived hireability judgement collected
from recordings; it is not an employer decision, a measure of job performance, or evidence that an
automated score is suitable for deployment. The original MIT Interview analysis introduced
automated prediction from mock interviews~\cite{Naim2015}, and its extended analysis examined
interview-performance prediction in more depth~\cite{Naim2018}. Comparisons of asynchronous with
face-to-face interviews raise a separate question about communication-skill measurement
~\cite{Rasipuram2016}; psychometric investigations of automated video-interview personality
assessments raise distinct reliability, validity and generalisation questions
~\cite{Hickman2022}.

The scientific question is whether a model trained without demographic attributes can use
regularisation across learned visual or audio groups to change disparity later observed across
audited groups, at a stated predictive cost. Answering it requires three distinct checks:
the underlying scorer must make useful predictions, the penalty must change the intended proxy
quantity on held-out people, and that change must be evaluated against separately sourced audit
attributes. A learned proxy is not a demographic identity. A visual judgement is not a person's
self-identified gender or ethnicity. The thesis therefore separates the intervention, the
held-out audit and the interpretation of each construct.

\section{Study corpora and empirical sequence}

The MIT Interview distribution contains 138 mock-interview sessions from 69 participants, each
recorded before and after coaching. Media, transcript and aggregate-outcome records join at the
source-session level. A corpus-wide identity check established the participant pairs and corrected
a reciprocal numbering error in the post-session recordings for participants 32 and 33 before
feature extraction or splitting. The observation is a session ($n=138$), while the independent
generalisation and uncertainty unit is a participant ($n=69$). Both sessions remain together in
every split. The primary target is the expectation--maximisation aggregate
\emph{RecommendHiring} rating.

The MIT representations use DINOv2~\cite{Oquab2024} for candidate-face crops and
WavLM-Base+~\cite{Chen2022WavLM} for bounded mixed-speaker interview audio from the same
active interval. MPNet supplies the text encoder's base~\cite{Song2020MPNet}; its sentence
representation follows Sentence Transformers~\cite{Reimers2019}, using the fixed
all-mpnet-base-v2 checkpoint~\cite{SentenceTransformersAllMPNet2025}.
Only audited candidate transcript turns enter the text representation. The audio does not
isolate the candidate and is not interpreted as a candidate-voice attribute.

The original neural experiment supplies a diagnostic starting point. It compares proxy W1
with an unchanged task-only model and proxy-matched sampling controls. A subsequent linear
follow-up uses the same cached representations and participant folds to examine a competitive
scorer, natural and balanced exposure, and the complete fixed penalty-strength grid. This
follow-up was motivated by the original audit, so its hypotheses are exploratory even though
perceived attributes remain excluded from fitting and selection.

AVI-Personality supplies an independent interview setting with six question-response clips
per participant and recruiter-rated outcomes~\cite{Zhang2026AVI}. The delivered release contains
646 participants and 644 available primary interview-performance ratings. Its official split,
participant-level video/audio features and self-reported gender source support a separate
prediction and error audit. Transcripts and skin-tone annotations are absent from the received
release. The different population, visual input boundary and rating construct prevent its raw
scores from being pooled with MIT's crowd rating or described as an identical replication.

\section{Research questions}

\begin{description}
    \item[RQ1] How predictable are the interview ratings for unseen participants after correcting
    source pairing, examining outcome measurement and comparing competitive simple scorers?

    \item[RQ2] Does increasing a proxy W1 penalty change held-out proxy-score disparity beyond
    the effect of balanced exposure, and what happens to predictive error and ranking?

    \item[RQ3] Do proxy changes transfer to phase-matched perceived-group score disparity or
    group errors, given single-observer annotation and sparse-category uncertainty?

    \item[RQ4] What prediction and group-error evidence is obtained in an independent interview
    corpus with different rating and attribute sources, and which conclusions remain
    specific to the MIT experiment?
\end{description}

These questions organise the thesis synthesis; they do not retrospectively make follow-up
analyses part of the original fixed design. Original contrasts remain separately reported.
The paired MIT recordings additionally support descriptive analysis of observed and predicted
post-minus-pre change. Without a randomised untreated comparison group, those differences
cannot identify a causal coaching effect.

\section{Study contributions}

The contributions form one connected argument. First, the thesis establishes a mock-interview
benchmark for unseen participants with an explicit pairing correction and outcome-reliability audit.
Second, it identifies the adverse incremental utility effect of adding W1 to matched sampling
controls in the original model. Third, a working linear follow-up tests the proxy mechanism
across the fixed strength grid and distinguishes held-out proxy response from uncertain transfer
to perceived-group disparities. Finally, independent video/audio prediction and self-reported-group
error results show how the inference changes with the corpus and rating construct. Separate
model and audit inputs support these comparisons as validity safeguards, rather than as the
headline research contribution.

The thesis does not claim that a learned proxy represents a protected class, that visual
annotations recover identity, or that distributional parity alone establishes hiring fairness.

\section{Thesis structure}

Chapter~\ref{ch:literature-review} reviews the relevant literature.
Chapter~\ref{ch:methodology} distinguishes the original design, exploratory linear follow-up and
independent-corpus analysis. Chapter~\ref{ch:results} reports prediction, proxy manipulation and
audited-group evidence in that order. Chapters~\ref{ch:discussion}--\ref{ch:conclusion} answer the
research questions and identify the limits of the conclusions.
